\section{What a limit tries to describe}

Limits are one of the central ideas of calculus. They are also one of the first ideas that require a new way of reading mathematical notation. A limit does not ask only what happens at a point. It asks what values are approached near a point, or far along a sequence.

This distinction is essential. The value of a function at a point and the limit of a function at that point are related, but they are not the same question.

\subsection*{Approaching is not the same as arriving}

Consider a function \(f\) and a point \(a\). The expression
\[
f(a)
\]
asks for the value of the function at the input \(a\). The expression
\[
\lim_{x\to a} f(x)
\]
asks what value \(f(x)\) approaches when \(x\) approaches \(a\).

The variable \(x\to a\) means that \(x\) gets close to \(a\). It does not mean that \(x=a\). This is why a limit can exist even when the function is not defined at the point.

\begin{examplebox}
Consider
\[
f(x)=\frac{x^2-1}{x-1},
\qquad x\ne1.
\]
The function is not defined at \(x=1\), because the denominator becomes zero. But for \(x\ne1\), we can factor:
\[
\frac{x^2-1}{x-1}=\frac{(x-1)(x+1)}{x-1}=x+1.
\]
When \(x\) approaches \(1\), the expression \(x+1\) approaches \(2\). Therefore
\[
\lim_{x\to1}\frac{x^2-1}{x-1}=2.
\]
The limit exists even though \(f(1)\) is not defined.
\end{examplebox}

This example should be read as a lesson about the question being asked. The limit does not ask what happens after substituting \(x=1\). It asks what values are approached when \(x\) gets close to \(1\).

\subsection*{Limits of sequences and limits of functions}

There are two main limit situations in this chapter.

For a sequence \((a_n)\), we ask what happens to the terms when the index grows:
\[
\lim_{n\to\infty} a_n.
\]
Here the input is discrete. The index \(n\) runs through natural numbers.

For a function \(f(x)\), we may ask what happens when \(x\) approaches a real number \(a\):
\[
\lim_{x\to a} f(x),
\]
or when \(x\) grows without bound:
\[
\lim_{x\to\infty} f(x).
\]

The notation is different because the input situation is different. But the common idea is the same: we study what values are approached.

\subsection*{What a limit can tell us}

Limits help us describe behavior that may not be visible from one single value. They allow us to discuss:
\begin{itemize}
\item long-term behavior of sequences,
\item behavior of functions near a point,
\item holes in graphs,
\item vertical and horizontal asymptotes,
\item continuity,
\item the idea of a derivative in the next chapter.
\end{itemize}

A limit is therefore not just another calculation. It is a way of reading behavior.

\begin{summarybox}
When first reading a limit, ask:
\begin{itemize}
\item What object is being studied: a sequence or a function?
\item What input is approaching something?
\item Is the input approaching a finite point or infinity?
\item Am I being asked about the value at the point or behavior near the point?
\item Does direct substitution answer the question, or does it reveal a difficulty?
\end{itemize}
\end{summarybox}

The main habit is this: before calculating a limit, read what kind of approaching is being described.